\origpage{352}{N. Robertson, P.D. Seymour / Journal of Combinatorial Theory, Series B 92 (2004) 325--357}
\hspace*{12pt}(8) \textit{For each $n \in N(\Gamma)$ there is a component $J$ of $H$ with $V(J) = V(\zeta(n))$; and for every component $J$ of $H$ not of this form with $E(J) \ne \emptyset$ there is a cuff $\Theta$ such that $V(J) \subseteq \bigcup (\beta(n) : n \in N(\Gamma) \cap \Theta)$.}

\textit{Subproof:} Every edge of $H$ either joins two nodes in $N(\Gamma')$ or joins two vertices both in $\bigcup (\beta(n) : n \in N(\Gamma) \cap \Theta)$ for some cuff $\Theta$. Let $n'_1, n'_2 \in N(\Gamma')$; we claim that they are connected in $H$ if and only if they both belong to $V(\zeta(n))$ for some $n \in N(\Gamma)$. First we prove the ``only if'' portion. If $n'_1, n'_2$ are connected in $H$ then they are joined by a path of $H$, all the vertices of which belong to $N(\Gamma')$, and so it suffices to prove the claim when $n'_1, n'_2$ are adjacent in $H$. Choose $c' \in C(\Gamma') \setminus \zeta(C(\Gamma))$ such that the edge of $H$ joining $n'_1, n'_2$ belongs to $\delta_{c'}$. If $c'$ is internal, then it follows from the definition of $\delta_{c'}$ that there exists $n \in N(\Gamma)$ with $n_1, n_2 \in V(\zeta(n))$ as required. If $c'$ is a border cell then from the seventh condition in the definition of ``linear inflation'', it follows that $n'_1, n'_2$ are adjacent in $V(\zeta(n))$ for some $n$, and again the claim holds. This proves ``only if''. Now for the ``if'' portion, assume that $n'_1, n'_2 \in V(\zeta(n))$. Since $\zeta(n)$ is a connected subgraph of $sk(\Gamma')$, we may assume that $n'_1, n'_2$ are adjacent in $sk(\Gamma')$ and hence in $V(\zeta(n))$. Let $c'$ be a cell of $\Gamma'$ such that $n'_1, n'_2 \in \tilde{c}'$. Since $\zeta(n)$ contains two different nodes of $c'$, it follows (from the sixth condition in the definition of ``linear inflation'') that $c' \notin \zeta(C(\Gamma))$. If $c'$ is internal, it follows that $n'_1, n'_2$ are adjacent in $H$ from the definition of $\delta_{c'}$, so we may assume that $c'$ is a border cell, and there is no internal cell $c'' \in C(\Gamma') \setminus \zeta(C(\Gamma))$ with $n'_1, n'_2 \in \tilde{c}''$. But then again it follows that $n'_1, n'_2$ are adjacent in $H$ from the definition of $\delta_{c'}$. This proves the ``if'' assertion, and thereby proves (8).

\hspace*{12pt}(9) \textit{Let $n \in N(\Gamma) \cap \Theta$, for some cuff $\Theta$. Let $n'_1, n'_2 \in V(\zeta(n)) \cap \Theta$ and let $1 \le i \le \rho(\Theta)$. Then the $i$th terms of $\pi'(n'_1)$ and $\pi'(n'_2)$ are connected in $H$.}

\textit{Subproof:} Since there is a line $F \subseteq \Theta$ such that for $n' \in N(\Gamma') \cap \Theta$, $n' \in V(\zeta(n))$ if and only if $n' \in F$, we may assume that $n'_1, n'_2$ are both nodes of some cell $c' \in C(\Gamma)$ bordering $\Theta$. Since $n'_1, n'_2 \in V(\zeta(n))$ it follows that $c' \notin \eta(C(\Gamma))$ and so $c' \ne c'_\Theta$ by (1). Hence $v'_1, v'_2$ are connected in $M'_{c'}$, from the defining property of $\pi'$, and hence they are connected in $H$. This proves (9).

\hspace*{12pt}(10) \textit{Let $n_1, n_2 \in N(\Gamma) \cap \Theta$ for some cuff $\Theta$, let $i > 0$, and let the $i$th term of $\pi(n_1)$ be the $i$th term of $\pi(n_2)$. Let $n'_1 \in V(\zeta(n_1)) \cap \Theta$ and $n'_2 \in V(\zeta(n_2)) \cap \Theta$. Then the $i$th term of $\pi'(n'_1)$ and the $i$th term of $\pi'(n'_2)$ are connected in $H$.}

\textit{Subproof:} By (9) the result holds if $n_1 = n_2$. Let $v$ be the $i$th term of $\pi(n_1)$. Since there is a line $F \subseteq \Theta$ such that for $n \in N(\Gamma) \cap \Theta$, $v \in \beta(n)$ if and only if $n \in F$, we may assume (by the argument used in the proof of (3)) that $n_1, n_2$ are both nodes of some cell $c$ bordering $\Theta$. By (9) we may replace $n'_1$ by any other element of $V(\zeta(n_1)) \cap \Theta$, for the result holds for the old element if and only if it holds for the new; and hence we may assume that $n'_1$ and similarly $n'_2$ are nodes of $c' = \eta(c)$. Since $\psi(c) \le \psi'(c')$ and the $i$th term of $\pi(n_1)$ is the $j$th term of $\pi(n_2)$ we deduce that the $i$th term of $\pi'(n'_1)$ is the $j$th term of $\pi'(n'_2)$. This proves (10).
